\chapter{الكيانات}
\label{ch:builder-entities}

\section{ماذا ستبني في هذا الفصل؟}
\label{sec:builder-entities-intro}

نوع السجلات، ويسميه كيان «الكيان»، هو أساس كل تطبيق تبنيه: فهو يحدد المعلومات التي تُحفظ، وطريقة عرضها، ومن يستطيع العمل عليها. ويشرح هذا الفصل بناء نوع السجلات من أوله إلى نشره؛ فتنشئه من صفحة «الكيانات»، ثم تضيف حقوله في «المنشئ»، وترتّب صفحته وقوائمه، وتضبط قواعده وصلاحياته، ثم تراجع تغييراتك وتنشرها.

ويُحفظ كل تعديل تجريه في «المنشئ» مسودةً لا يراها أحد غيرك، فتستطيع أن تجرّب وتصحّح بأمان. ولا تصل التغييرات إلى المستخدمين إلا بعد نشرها، لذلك تنتهي كل مرحلة في هذا الفصل بالمراجعة والنشر.

\begin{rcnote}[مثال: شركة الأفق]
يبني مسؤول النظام في الأفق نوع السجلات «طلبات الزيارة»: يضيف حقوله، ويرتّب صفحته، ويسمح لعمر حداد بإضافة الطلبات ولسارة الخطيب بمراجعتها، ثم ينشره ليبدأ الفريق باستخدامه.
\end{rcnote}

\begin{figure}[H]
\centering
\begin{tikzpicture}[node distance=0.5cm]
\node[entity] (type) {\RL{نوع السجلات}};
\node[entity,below=0.7cm of type] (fields) {\RL{الحقول}};
\node[entityOpt,below right=0.7cm and 0.3cm of fields] (page) {\RL{صفحة السجل}};
\node[entityOpt,below left=0.7cm and 0.3cm of fields] (lists) {\RL{القوائم وطرق العرض}};
\draw[dependsOn] (type) -- (fields);
\draw[dependsOnOpt] (fields.south) -- ++(0,-0.3) -| (page.north);
\draw[dependsOnOpt] (fields.south) -- ++(0,-0.3) -| (lists.north);
\end{tikzpicture}
\rclegend
\caption{أجزاء نوع السجلات}
\label{fig:builder-models}
\shotnote{يوضّح الشكل أن الحقول جزء لازم من نوع السجلات، لأنها تحدد المعلومات التي يحفظها كل سجل. أما صفحة السجل والقوائم فطريقتان لعرض هذه الحقول، ولكل منهما إعداد افتراضي يعمل دون تعديل، فتضبطهما عندما تريد عرضاً أوضح للفريق.}
\end{figure}

\section{صفحات هذا الفصل}
\label{sec:builder-entities-pages}

تبدأ الصفحتان الأوليان من قائمة «الكيانات» في قسم «الأساسية»، أما بقية الصفحات فتشرح أقسام داخل «المنشئ» تصل إليها بفتح نوع السجلات. وترتيب البطاقات هو ترتيب البناء المقترح، فاتبعه في المرة الأولى، ثم عُد إلى أي صفحة عندما تحتاج إليها.

\begin{table}[H]
\centering
\caption{بناء نوع السجلات خطوة بخطوة}
\label{tab:builder-pages-entities}
\rcZebra
\begin{tabularx}{\linewidth}{R{4.2cm} Y}
\rcHeaderRow
\thh{الصفحة وقسمها} & \thh{ماذا تقدم هذه الصفحة؟}\\
الكيانات — الأساسية & تعرض أنواع السجلات في المؤسسة، ومنها تنشئ نوعاً جديداً أو تفتحه للتعديل.\\
منشئ الكيانات بالذكاء الاصطناعي ومخطط العلاقات — الأساسية & تصف النوع بكلماتك فيقترح المساعد حقوله، وترى العلاقات بين الأنواع قبل أن تعدّلها.\\
المنشئ والنظرة العامة — المنشئ & رأس المنشئ وأقسامه، وبطاقات «نظرة عامة» وقائمة «قبل النشر».\\
الحقول — المنشئ & تحدد الحقول المعلومات التي يحفظها كل سجل ونوع كل معلومة.\\
أنواع الحقول وإعداداتها — المنشئ & كل نوع حقل ومتى تستخدمه، والإعدادات الإضافية التي يطلبها بعض الأنواع.\\
مصدر القيمة: الحقول المحسوبة والمرتبطة والمجمّعة — المنشئ & قيم تُحسب بصيغة أو تُنقل من سجل مرتبط أو تُجمع من السجلات المرتبطة، ومتى يكون الحقل إلزامياً.\\
حقل الحالة وقواعد دورة الحياة — المنشئ & شروط الانتقال بين الحالات، ومن لا يحق له تنفيذه، وما يُكتب عند الانتقال، وقفل السجل.\\
صفحة السجل — المنشئ & ترتّب الصفحة التي يراها المستخدم عند فتح السجل في تبويبات وأقسام.\\
كتل النظرة العامة والتخطيطات — المنشئ & ملخص أعلى صفحة السجل بالمؤشرات والتنبيهات، وتخطيطات تختلف حسب قيمة في السجل.\\
القوائم وطرق العرض — المنشئ & تحدد كيف تظهر السجلات معاً: جدولاً أو لوحة أو تقويماً، وما الأعمدة الظاهرة.\\
المؤشرات والعروض المحفوظة — المنشئ & أرقام تلخيصية فوق القائمة، وتبويبات بتصفية وترتيب جاهزين.\\
القواعد — المنشئ & تُظهر الحقول أو تقفلها أو تُلزم بها أو تتحقق من قيمها بحسب ما في السجل.\\
الحدود وحراس المراجع والتوليد والنطاقات الزمنية — المنشئ & ضوابط بين السجلات تمنع تجاوز المبالغ وحذف السجلات المستخدمة وتعارض التواريخ.\\
الصلاحيات — المنشئ & تحدد الأدوار التي تستطيع رؤية السجلات وإضافتها وتغييرها وحذفها.\\
صلاحيات الحقول والتحقق من صلاحيات شخص — المنشئ & من يرى كل حقل ومن يغيّره، وفحص ما يراه شخص معيّن قبل النشر.\\
الإجراءات والأتمتة — المنشئ & تضيف إلى السجل أزراراً تُنشئ مستندات أو تفتح روابط أو تبدأ مسارات عمل.\\
الحركة — المنشئ & تنقل مبلغاً بين سجلات نوع آخر وترحّله عند حالة معينة، للأنواع المالية والمخزنية.\\
الاستيراد — المنشئ & يحدد كيف تُحمَّل السجلات وتُحدَّث من ملفات Excel، وبأي مفتاح تُطابق الصفوف مع السجلات.\\
إعدادات متقدمة — المنشئ & اسم السجل ومنع التكرار والسجلات الأم ومكان ظهور النوع.\\
المراجعة والنشر — المنشئ & تراجع التغييرات المحفوظة كمسودة، ثم تنشرها ليراها الفريق، وتعود إلى إصدار سابق عند الحاجة.\\
\end{tabularx}
\end{table}
